\chapter{El amor no es un cuento}

\epigraph{Amar no es mirarse el uno al otro, sino mirar juntos en la misma dirección.}{Antoine de Saint-Exupéry}

Todos hemos visto la escena. Él y ella se miran a los ojos. Suena la música. El mundo se detiene. Se besan. Aparecen las letras: ``Y vivieron felices para siempre''. Fin.

La película acaba justo donde empieza lo interesante.

Porque si algo hemos aprendido en los capítulos anteriores, es que nuestro cuerpo está extraordinariamente preparado para el encuentro. Tenemos hormonas que nos empujan a vincularnos, circuitos neuronales que nos hacen sentir euforia junto a la persona que nos atrae, mecanismos biológicos que favorecen la permanencia. Todo eso es real, está inscrito en nuestra naturaleza, y es asombroso.

Pero hay un salto enorme entre sentir todo eso y construir un amor que dure. Un salto que las películas se saltan alegremente ---y que la vida, con su habitual falta de banda sonora, nos obliga a dar con los pies en el suelo.

\section{El mito del amor de película}

No es que las películas de amor sean malas. Algunas son estupendas. El problema es cuando convertimos la ficción en modelo de lo que el amor debería ser. Y eso, sin darnos cuenta, lo hacemos más de lo que pensamos.

El guion suele ser parecido. Dos personas se encuentran. Hay una chispa inmediata, un flechazo que ninguno de los dos puede controlar. Quizá hay un obstáculo ---un malentendido, una distancia, un rival--- que se resuelve en el tercer acto con una carrera bajo la lluvia o una declaración épica en un aeropuerto. Y luego: la felicidad. El amor verdadero ha triunfado.

¿Qué nos enseña ese guion, si lo miramos con un poco de distancia?

Primero, que el amor es algo que te \textit{pasa}. Que llega sin que lo busques, te arrolla como un tren y ya está. Si no sientes mariposas en el estómago, si no pierdes el sueño, si no te tiemblan las piernas, entonces no es ``el de verdad''. El amor auténtico se reconoce porque es irresistible.

Segundo, que el amor no debería costar esfuerzo. Si hay que trabajar en la relación, si hay que pedir perdón, si hay que negociar quién friega los platos, algo falla. El amor de verdad fluye. Es natural. Es fácil.

Y tercero ---quizá lo más dañino---, que existe una única persona perfecta para ti en algún lugar del mundo. Tu ``media naranja'', tu alma gemela, la pieza que encaja sin necesidad de limar bordes. Y si un día descubres que la persona que tienes al lado no es exactamente esa pieza... bueno, quizá te equivocaste. Quizá hay que seguir buscando.\footnote{El sociólogo Zygmunt Bauman analiza con lucidez esta mentalidad consumista aplicada a las relaciones en \textit{Amor líquido}, Fondo de Cultura Económica, 2005. Cuando el otro deja de satisfacernos, lo ``devolvemos'' y buscamos un modelo mejor.}

Dicho así, suena absurdo. Pero conviene ser honestos: ¿cuántas relaciones se han roto no porque hubiera un problema grave, sino porque alguien esperaba sentir permanentemente lo que sintió los primeros meses? ¿Cuántas personas han abandonado algo bueno ---algo real--- persiguiendo una intensidad que, por definición, no puede durar para siempre?

\section{Cuando la química se confunde con el amor}

Aquí es donde lo que vimos en el capítulo anterior cobra una importancia decisiva. Recordemos: el enamoramiento tiene un sustrato biológico poderosísimo. Dopamina a raudales, noradrenalina que nos mantiene en alerta, serotonina baja que nos hace pensar obsesivamente en la otra persona. Es un cóctel químico diseñado por la evolución para que dos personas se acerquen, se deseen, quieran estar juntas. Y funciona. Vaya si funciona.

Pero tiene fecha de caducidad.

No es que el sentimiento desaparezca del todo. Es que cambia. La intensidad del inicio ---ese no poder pensar en otra cosa, esa euforia, ese ``me muero si no te veo''--- se atenúa con el tiempo. Y tiene que ser así. Ningún organismo puede mantenerse indefinidamente en ese estado de excitación sin colapsarse. Los estudios sugieren que la fase de enamoramiento intenso dura, en la mayoría de los casos, entre uno y tres años.\footnote{Cf. Fisher, H., \textit{Why We Love: The Nature and Chemistry of Romantic Love}, Henry Holt, 2004. Fisher distingue tres sistemas cerebrales del amor ---impulso sexual, atracción romántica y apego--- con sustratos neuroquímicos diferenciados.}

Y cuando baja la marea ---cuando la dopamina deja de inundar cada encuentro, cuando puedes volver a concentrarte en el trabajo, cuando duermes otra vez de un tirón--- llega el momento de la verdad. Porque ahí es donde muchas personas sienten que algo se ha roto. ``Ya no siento lo mismo'', dicen. ``Se ha acabado la magia.''

Pero lo que se ha acabado no es el amor. Lo que se ha acabado es la fase química del enamoramiento. Son cosas distintas. Tan distintas como el fogonazo que enciende la chimenea y el fuego lento que calienta la casa durante todo el invierno. El fogonazo es espectacular, sí. Pero es el fuego sostenido el que te salva del frío.

Confundir el enamoramiento con el amor es como confundir el hambre con la nutrición. El hambre te lleva a la mesa. Pero lo que te alimenta es comer ---y comer bien, día tras día, no un atracón puntual seguido de ayuno.

\section{Amar es un verbo}

Erich Fromm escribió hace más de medio siglo algo que sigue siendo revolucionario por lo sencillo: el amor no es un sustantivo, es un verbo. No es algo que se \textit{tiene}, sino algo que se \textit{hace}.\footnote{Fromm, E., \textit{El arte de amar}, Paidós, 1956. Fromm argumenta que el amor es un arte que requiere conocimiento, esfuerzo y práctica, igual que cualquier otra actividad humana seria.}

Esto cambia todo.

Si el amor fuera solo un sentimiento ---algo que te pasa, que te invade, que no controlas---, entonces no tendría mucho sentido prometer amor. Sería como prometer que mañana va a hacer sol. Puedes desearlo, pero no depende de ti.

Sin embargo, si el amor es también ---y sobre todo--- una decisión, una acción, algo que haces con tu voluntad libre, entonces sí puedes prometerlo. Porque puedes decidir cuidar al otro aunque hoy no sientas mariposas. Puedes elegir escuchar cuando estás cansado. Puedes optar por el perdón cuando sería más fácil llevar la cuenta de los agravios. Puedes levantarte cada mañana y renovar la decisión de querer a esa persona concreta, con sus virtudes y con sus manías, con su luz y con sus sombras.

Eso no significa que el sentimiento no importe. Claro que importa. El deseo, la ternura, la emoción de estar con alguien que te hace bien: todo eso es parte del amor y sería triste prescindir de ello. Pero el sentimiento solo no basta. Sin la decisión, el sentimiento es como un barco sin timón: va donde lo lleve el viento.

Fijaos en algo curioso. Cuando un padre o una madre dicen ``quiero a mi hijo'', nadie entiende que eso signifique que sienten una emoción agradable cada vez que lo miran. Hay noches en que el niño llora y tú estás agotado. Hay adolescencias que ponen a prueba la paciencia de un santo. Hay momentos en que el amor parental se parece más a la determinación que a la ternura. Y sin embargo, nadie duda de que eso sea amor. Lo es, y del bueno.

¿Por qué, entonces, con el amor de pareja exigimos una montaña rusa emocional permanente? ¿Por qué si un día el corazón no se acelera pensamos que algo va mal?

\section{Las tres caras del amor}

Los griegos, que tenían una palabra para casi todo, usaban al menos tres términos distintos donde nosotros usamos uno solo: ``amor''.\footnote{La distinción es más amplia aún. El griego contaba también con \textit{storgé} (amor familiar, afecto natural) y otros matices. Aquí nos centramos en las tres dimensiones más relevantes para el amor de pareja.} Y esa distinción no es un capricho filológico: ilumina algo que la experiencia confirma.

\subsection*{Eros: el deseo que busca}

El eros es la primera chispa. Es la atracción, el deseo, la pasión. Es esa fuerza que te lleva a querer estar cerca del otro, a tocarlo, a mirarlo. Es el componente del amor que siente el cuerpo antes de que la cabeza pueda explicar lo que pasa.

El eros es bueno. No es algo de lo que avergonzarse ni algo que haya que domesticar como a una bestia. Es parte de nuestra naturaleza, y ya hemos visto que tiene raíces biológicas profundas. Sin eros, probablemente dos personas nunca se acercarían lo suficiente como para descubrir lo que tienen en común.

Pero el eros tiene una limitación. Es, por naturaleza, un amor que busca para sí. Deseo al otro porque me hace sentir bien, porque me completa, porque me hace feliz. No hay nada malo en eso ---¿quién no quiere ser feliz?---, pero si el amor se queda ahí, se vuelve frágil. Porque el día que el otro deje de hacerme sentir bien ---porque está enfermo, o cansado, o simplemente porque los años han pasado y la novedad se ha agotado---, ¿qué queda?

\subsection*{Philia: la amistad que comparte}

El segundo rostro del amor es lo que los griegos llamaban \textit{philia}: la amistad profunda. Es el amor que nace de la vida compartida, de los intereses comunes, del placer de estar juntos no solo por la atracción, sino por la compañía.\footnote{C.\,S. Lewis dedica un capítulo espléndido a la amistad en \textit{Los cuatro amores} (1960). Para Lewis, la amistad es el menos biológico de los amores y, quizá por eso, el más cercano a lo angélico: ``Los amigos se miran el uno al otro; los enamorados se miran a los ojos.''}

Hay parejas que se desean pero no se gustan. Hay parejas con una química explosiva que, si les quitas la cama, no tienen de qué hablar. El eros sin philia es un fuego que consume pero no calienta. La philia es lo que convierte a la persona que amas en tu amigo, tu confidente, tu compañero de viaje. La persona con la que puedes estar en silencio sin que el silencio pese.

Pensadlo un momento: en un matrimonio que dure toda la vida, ¿cuántas horas se pasan conversando, comiendo juntos, decidiendo cosas prácticas, riendo, discutiendo, resolviendo problemas, sentados en el mismo sofá? Muchas más que en cualquier otra actividad. Si esa persona no es también tu amigo, tienes un problema.

La philia crece donde el eros la planta. Pero necesita su propio cuidado: tiempo compartido, intereses cultivados en común, conversación honesta. Parejas que dejan de hablar ---de hablar de verdad, no de logística doméstica--- están dejando morir la philia sin darse cuenta.

\subsection*{Ágape: el amor que se entrega}

Hay un tercer nivel, y es el más difícil. El más exigente. Y, paradójicamente, el que hace que el amor sea verdaderamente grande.

El ágape es el amor como don. No el amor que busca (eros) ni el amor que comparte (philia), sino el amor que se da sin condiciones. Es el amor que dice: ``Quiero tu bien, aunque me cueste. Aunque no me lo agradezcas. Aunque hoy no me apetezca.''

Es la madre que se levanta a las tres de la mañana a atender al bebé. Es el marido que acompaña a su mujer a la quimioterapia sin faltar un día. Es la esposa que sostiene a su marido cuando ha perdido el trabajo y la confianza en sí mismo. No porque sea agradable hacerlo, sino porque el otro lo necesita y yo puedo darlo.

Fromm lo expresó con una fórmula que merece la pena recordar: el amor inmaduro dice ``te amo porque te necesito''; el amor maduro dice ``te necesito porque te amo''.\footnote{Fromm, \textit{El arte de amar}, cap. II. La fórmula es sencilla pero profundamente reveladora: en el amor maduro, la necesidad nace del amor y no al revés.} La diferencia es abismal. En el primer caso, el otro es un medio para satisfacer mi necesidad. En el segundo, mi necesidad nace de haber descubierto en el otro un valor que merece todo lo que puedo dar.

El ágape no anula al eros ni a la philia. Los necesita. Un matrimonio solo de ágape, sin deseo ni amistad, sería sacrificio heroico pero no plenitud. Las tres dimensiones se necesitan como las patas de un trípode: quita una y todo se cae.

Pero es el ágape el que sostiene al amor cuando las otras dos dimensiones pasan por sus inevitables valles. Y todos los matrimonios tienen valles. Es el ágape lo que te lleva a elegir al otro cuando elegirlo cuesta. Y es ahí, precisamente ahí, donde el amor revela su verdadero rostro.

\section{Conocerte para poder darte}

Hay una verdad incómoda que conviene mirar de frente: no se puede dar lo que no se tiene.

Si no me conozco ---si no sé cuáles son mis heridas, mis miedos, mis necesidades profundas, mis mecanismos de defensa---, lo que llamaré ``amor'' será muchas veces otra cosa disfrazada. Será dependencia emocional vestida de entrega. Será miedo a la soledad vestido de compromiso. Será la búsqueda inconsciente de que el otro repare lo que yo no he querido mirar en mí mismo.

Esto no es un reproche. Es una realidad que conviene nombrar antes de prometer algo tan grande como ``te amaré toda la vida''. Porque esa promesa exige un sujeto que sepa quién es y qué está prometiendo.

La madurez afectiva no es un estado al que se llega un buen día y ya está. Es un camino. Pero sí hay señales que indican si estamos en ese camino o si estamos evitándolo. Quien se conoce, sabe distinguir entre lo que necesita y lo que desea. Sabe poner nombre a sus emociones sin que lo arrastren. Es capaz de estar solo sin sentir pánico ---y precisamente por eso es capaz de estar con otro sin asfixiarlo---.\footnote{La psicóloga Harriet Lerner lo resume con claridad: ``Solo puedes ser realmente íntimo con otro si no te has perdido a ti mismo.'' Cf. Lerner, H., \textit{La danza de la intimidad}, Urano, 1990.}

Quien no se conoce, en cambio, tiende a buscar en la pareja lo que debería buscar en sí mismo. Y eso es cargar sobre los hombros del otro un peso que no le corresponde. Nadie puede ser tu terapeuta, tu padre ausente, tu autoestima y tu pareja al mismo tiempo. Es demasiado. Y antes o después, la estructura se resquebraja.

Decimos muchas veces que el matrimonio es ``darse al otro''. Y es verdad. Pero solo puedes darte si te tienes. Solo puedes hacer un regalo de ti mismo si sabes qué hay en el paquete.

\section{La paradoja de la libertad}

Hay una objeción que aparece una y otra vez, especialmente entre los más jóvenes: ``¿Para qué comprometerse? ¿No es mejor mantener las opciones abiertas?''

La lógica parece impecable. Si me comprometo con una persona, me cierro a todas las demás. Si no me comprometo, puedo elegir cada día lo que más me convenga. Máxima libertad, máxima flexibilidad.

Pero la experiencia dice otra cosa. Y conviene escucharla.

Pensemos en alguien que, por miedo a equivocarse, nunca elige una carrera profesional. Va probando un poco de esto, un poco de aquello. Mantiene ``las opciones abiertas''. A los cuarenta años, tiene mucha experiencia superficial y ninguna maestría. No ha profundizado en nada. No es el profesional brillante que podría haber sido, porque serlo habría requerido años de dedicación a una sola cosa.

Con el amor pasa algo parecido. El compromiso no es la tumba de la libertad; es su cauce. Como un río: el agua que se desborda por todas partes no llega a ningún sitio; el agua que fluye entre riberas puede mover molinos, regar campos, llegar al mar.

Cuando eliges a una persona ---cuando dices ``tú, y no otra''---, no estás limitando tu capacidad de amar. La estás concentrando. La estás dirigiendo. Y al hacerlo, puedes llegar a una profundidad que el que va saltando de relación en relación nunca conocerá.

Hay algo más, y es quizá lo más sorprendente. El compromiso genera una forma de libertad que la indecisión no puede dar: la libertad de mostrarte como eres.

Pensad en la diferencia entre una primera cita y una cena con alguien que llevas veinte años queriendo. En la primera cita te muestras en tu mejor versión. Elegiste la camisa perfecta, preparaste tres anécdotas graciosas, te cuidas de no decir nada inconveniente. Estás representando un papel.

Después de veinte años, no. Después de veinte años puedes estar en pijama, con la cara lavada, diciendo tonterías, mostrando tus inseguridades, y sabes que el otro no se va a ir. Esa seguridad ---la certeza de que el otro ha elegido quedarse y va a seguir eligiéndolo--- es lo que te permite dejar de actuar y empezar a vivir.

El filósofo Gabriel Marcel decía que la fidelidad creadora no es la repetición monótona de una decisión pasada, sino la renovación constante de un compromiso vivo.\footnote{Cf. Marcel, G., \textit{Être et avoir}, Aubier, 1935. Marcel distingue entre la fidelidad inerte ---la del que se queda por inercia--- y la fidelidad creadora, que reinventa cada día las razones para quedarse.} No te quedas porque lo prometiste hace años y ``no queda otra''. Te quedas porque cada día descubres nuevas razones para elegir lo que ya elegiste. Y eso ---elegir libremente lo que ya es tuyo--- es una de las formas más puras de la libertad.

\section{Un amor que crece}

Terminemos volviendo al principio. El cuento dice que el amor verdadero es el del flechazo, el del corazón desbocado, el del ``lo supe desde el primer momento''. Y que si no es así, no es de verdad.

La realidad dice otra cosa. La realidad dice que los amores más hondos, los más sólidos, los que dan más calor, son los que se han construido ladrillo a ladrillo. Los que han pasado por crisis y han salido más fuertes. Los que han aprendido a pedir perdón y a concederlo. Los que se han reinventado cuando la vida les cambió las reglas del juego.

Hay una belleza particular en las parejas que llevan décadas juntas. No la belleza de la portada de revista, sino la belleza de algo que ha sido probado por el tiempo y ha resistido. La belleza de dos personas que se miran y, debajo de las arrugas y las canas, siguen viendo al otro que eligieron ---y lo siguen eligiendo hoy.

Eso no lo da un fogonazo. Lo da un fuego que alguien ha tenido la paciencia y el coraje de alimentar, día tras día, año tras año.

El amor no es un cuento. Es algo mucho mejor que un cuento. Porque los cuentos terminan, y el amor verdadero no tiene por qué hacerlo.

\paracaminar

\begin{enumerate}
  \item ¿Reconocéis en vuestra propia historia algún ``mito'' sobre el amor que hayáis tenido que desmontar? ¿Qué esperabais del amor antes de vivirlo y qué habéis descubierto después?

  \item De las tres dimensiones del amor ---eros, philia, ágape---, ¿cuál sentís que está más presente en vuestra relación ahora mismo? ¿Cuál necesita más cuidado?

  \item ¿En qué momento concreto habéis experimentado que elegir al otro ---cuando costaba, cuando no apetecía--- ha hecho crecer vuestro amor en vez de reducirlo?
\end{enumerate}

\vinetacapitulo{ch04}
